


\label{block:full}


\subsection{Objetivo}
Cuantificar el impacto de la remodelación de estaciones de servicio al formato \textbf{YPF FULL} —estaciones con tienda de conveniencia modernizada y servicios ampliados— sobre las ventas de combustibles y de la tienda asociada. En particular, se busca estimar en qué
medida la conversión al formato FULL incrementa el volumen de combustible despachado, la facturación total y las ventas de conveniencia, en relación con la trayectoria contrafactual de estaciones comparables que no realizaron dicha inversión. Este grupo de comparación puede incluir tanto estaciones YPF que permanecen en formato tradicional como estaciones de la competencia con características similares (por ejemplo, ubicación, tamaño y tendencia previa de ventas) que no implementan mejoras sustantivas en la calidad de sus servicios durante el período de análisis. 

Adicionalmente, el estudio puede extenderse a la evaluación del beneficio económico neto de la conversión, incorporando información sobre los costos de implementación del formato FULL, tanto de inversión inicial (costos fijos de remodelación) como de operación (costos variables y operativos), lo que permitiría construir métricas de retorno de la inversión y contribuir a la priorización de futuras decisiones de expansión del formato, así como al análisis de la creación de nuevos espacios premium dentro de las estaciones, tales como iniciativas recientemente anunciadas bajo la marca \textbf{YPF Black}.

%\subsection{Justificación y relevancia} La inversión requerida para la remodelación de estaciones de servicio es significativa, por lo que resulta fundamental cuantificar su retorno económico en términos de mayores ventas y márgenes. En particular, las tiendas \emph{Full} se han consolidado como un componente central del negocio, llegando a representar una fracción sustantiva de la rentabilidad de una estación (\textcolor{red}{chequear}). En línea con esta tendencia, YPF ha registrado un fuerte crecimiento en el desempeño de sus tiendas de conveniencia —incluyendo una duplicación de las ganancias en un año y la expansión hasta aproximadamente 1.000 tiendas Full hacia fines de 2024 (\textcolor{red}{chequear})— lo que subraya la relevancia estratégica de este formato.

% Más allá del aumento agregado de ventas, resulta particularmente relevante evaluar si la conversión al formato Full induce una mayor demanda de combustibles \emph{premium}, que presentan márgenes superiores, y si dichos consumidores exhiben patrones de consumo complementarios dentro de la tienda de conveniencia, caracterizados por la compra de productos de mayor calidad y mayor margen. Este canal sugiere un potencial efecto multiplicador de las inversiones en formato Full, al reorientar la demanda hacia segmentos más rentables tanto en combustibles como en productos de la tienda.

% La evidencia disponible para el mercado local respalda la hipótesis de un efecto positivo de las mejoras en servicios complementarios: por ejemplo, se ha documentado que la introducción de estándares gastronómicos en estaciones de servicio incrementó las ventas de combustibles en torno al 10\% en promedio, con aumentos del orden del 20--25\% en las ventas de naftas premium (Cerutti et al., 2026). En este contexto, contar con estimaciones precisas y causalmente identificadas del impacto de la transformación a Full resulta clave para la toma de decisiones. Diferencias en magnitudes aparentemente moderadas —por ejemplo, un aumento del 6\% frente a uno del 9\%— implican variaciones económicamente muy relevantes cuando se las evalúa sobre volúmenes elevados de ventas de combustibles. Por ello, la obtención de coeficientes robustos y bien identificados es esencial para evaluar el retorno de las remodelaciones y optimizar futuras decisiones de inversión, como la priorización de ubicaciones a remodelar o la expansión de formatos premium dentro de la red.

\subsection{Diseño econométrico} Se propone un diseño de \emph{diferencias en diferencias} con tratamiento escalonado (\emph{staggered adoption}) implementado en un marco de \emph{estudio de eventos}, es decir, un único esquema de estimación que combina ambas herramientas siguiendo metodologías de frontera recientes para tratamientos que se adoptan en distintos momentos. En términos prácticos, se compararán las trayectorias de ventas de combustibles y de tienda de las estaciones que se convierten a Full (grupo de tratamiento) antes y después de la remodelación, frente a estaciones que aún no han pasado al formato Full en el mismo período (grupos de control). Este enfoque permite, por un lado, verificar si antes de la intervención los grupos comparados seguían trayectorias similares —una condición necesaria para que la comparación sea confiable— y, por otro, estimar cómo evoluciona el impacto en los períodos posteriores, evitando los sesgos que pueden surgir en métodos de comparación más simples.

Para enriquecer el análisis, se explorarán mecanismos específicos. Por ejemplo, se distinguirá si el aumento en las ventas de combustibles proviene de una mayor afluencia de vehículos (\emph{tráfico}) o de un mayor consumo por cliente (litros por transacción), así como el impacto sobre la mezcla de productos (incrementos relativamente mayores en combustibles premium respecto de los regulares, lo que sería consistente con la atracción de clientes de mayor disposición a pagar). De manera análoga, si se dispone de información pre y post conversión a Full sobre las cantidades vendidas en la tienda de conveniencia —por categoría de producto o incluso a nivel de ítem— se podrá analizar si el formato Full modifica la composición del consumo en la tienda (por ejemplo, hacia productos de mayor calida y margen unitario), complementando así la evidencia sobre el canal de servicios adicionales.

Adicionalmente, se podrá variar sistemáticamente la definición de los grupos de control —considerando distintos conjuntos de estaciones YPF comparables y, cuando los datos lo permitan, estaciones de la competencia similares en ubicación y tendencia previa— con el objetivo de evaluar la robustez de los resultados frente a diferentes contrafactuales plausibles.

\subsection{Datos requeridos} Datos a nivel de estación de servicio con periodicidad mensual (o, preferiblemente, diarios si estuvieran disponibles) que incluyan el volumen de combustibles vendidos —desagregado por tipo de combustible (por ejemplo, nafta súper, nafta premium, gasoil grado 2 y gasoil grado 3)— y los ingresos totales (mensuales o diarios) de las tiendas de conveniencia para cada estación. Para el análisis básico de combustibles, esta información puede obtenerse a partir de los datos administrativos disponibles en la Resolución del Ministerio de Energía, que permiten construir paneles de ventas por estación y mes \textcolor{red}{Alguna oración que diga que es preferible con sus datos porque los del ministerio tienen muchos errores de carga}. Asimismo, se requiere el listado de estaciones con la fecha de remodelación y conversión al formato Full (historial de implementación). \textcolor{red}{esto ya fue provisto por YPF}

Para análisis más detallados y de mayor profundidad, resulta deseable complementar esta información con datos operativos internos de YPF. En particular, contar con variables a nivel de transacción o, alternativamente, con información agregada sobre el número de transacciones y la cantidad total de litros vendidos por período permitiría distinguir entre aumentos en el tráfico de clientes y cambios en el consumo promedio por visita (litros por carga). De manera análoga, para la tienda de conveniencia sería valioso disponer de información sobre cantidades vendidas por categoría de producto, número de tickets y ticket promedio, lo que permitiría analizar cambios en la composición del consumo y en la venta de productos de mayor margen.

Adicionalmente, se podrán incorporar variables de control relativas a las características de la estación (capacidad instalada, número de surtidores, ubicación urbana o sobre rutas) y, cuando sea posible, indicadores del entorno económico o de tráfico vehicular. En caso de extender el análisis a la dimensión competitiva, sería útil contar también con información sobre ventas de combustibles de estaciones de la competencia en la misma zona, ya sea a partir de fuentes públicas, datos administrativos o estimaciones de mercado disponibles.




